% ===========================================================================
%  HALO — IEEE TIP Draft
%  Introduction + Related Work + Methodology
%  Last updated: 2026-03-17 (Sec.2.3 expanded)
% ===========================================================================

% ===========================================================================
\section{Introduction}\label{sec:intro}
% ===========================================================================

Infrared small target detection (IRSTD) plays an indispensable role in maritime
surveillance, early-warning systems, and remote sensing~\cite{dnanet,uiunet,alcnet}.
Unlike generic object detection, infrared small targets typically occupy fewer than
$9\times9$ pixels, lack discernible texture or shape cues, and are deeply embedded
in spatially correlated background clutter~\cite{irstd_survey}.
Despite the remarkable progress brought by deep learning~\cite{dnanet,uiunet,sctransnet},
fully supervised methods remain fundamentally bottlenecked by the cost of acquiring
pixel-level mask annotations.
Precisely delineating a $3\times3$ target against a noisy infrared background demands
domain expertise and meticulous per-pixel labeling---a process that scales poorly to
the large-scale datasets required for robust generalization.

To mitigate this annotation burden, recent works have explored \emph{weakly supervised}
paradigms for IRSTD.
Notably, LESPS~\cite{lesps} introduces single-point supervision with a label evolution
mechanism that iteratively refines pseudo-masks during training, while PAL~\cite{pal}
further incorporates an active learning loop to progressively select the most informative
annotation points.
Although these methods significantly reduce annotation effort, they share two fundamental
limitations that hinder deployment in operational scenarios.
\emph{First}, single-point annotation in dense clutter is inherently fragile: accurately
clicking the geometric centroid of a dim, sub-$5$-pixel target embedded in heavy
sea-surface or cloud clutter is both time-consuming and prone to human jitter, as
illustrated in Fig.~\ref{fig:annotation_comparison}.
\emph{Second}, the iterative label evolution and active learning loops employed by LESPS
and PAL introduce substantial computational overhead and training complexity, undermining
the simplicity required for real-world deployment.

In contrast, \emph{bounding box} annotations represent the most natural and readily
available form of weak supervision in practice.
Upstream early-warning systems---ranging from radar cueing to lightweight YOLO-based
detectors~\cite{yolo}---routinely produce coarse bounding boxes as their primary output.
Requiring these operational pipelines to additionally supply pixel-level masks or
geometrically precise center points is neither realistic nor cost-effective.
However, leveraging box annotations for IRSTD poses a unique and severe challenge:
for a typical small target occupying fewer than $0.5\%$ of the box area, over $99\%$
of the enclosed pixels belong to background clutter.
Naively treating all box-interior pixels as foreground---so-called \emph{Box Fill}---leads
to catastrophic \emph{representation collapse}, where the network learns to predict the
entire box region rather than the compact target, yielding performance \emph{below}
single-point supervised baselines.
This observation is empirically confirmed in our experiments: Box Fill (55.17\%,
57.39\%, 50.90\% mIoU on NUAA/NUDT/IRSTD) and GrabCut (57.31\%, 42.10\%, 50.43\%)
both fall short of LESPS (61.13\%, 58.37\%, 49.05\%), demonstrating that the bounding
box itself carries no inherent information advantage---the challenge lies precisely
in extracting the target signal from within it.

In this paper, we propose \textbf{HALO} (\textbf{H}eat-\textbf{A}nchored \textbf{L}abel
\textbf{O}ptimization from Loose Boxes)%
\footnote{The name also alludes to the faint optical ``halo'' surrounding an infrared
point source, a direct physical manifestation of the PSF exploited by our method.},
a physics-anchored pseudo-label generation framework that converts coarse bounding boxes
into high-fidelity pixel-level soft masks \emph{without any learnable parameters,
iterative refinement, or network architecture modification}.
Our key insight is that the radiometric properties of infrared imagery provide a
principled statistical basis for distinguishing target pixels from background within
a bounding box.
Concretely, we contribute two components at the pseudo-label generation level:

\begin{enumerate}
    \item \textbf{Background-SNR Posterior Mask (B-SNR).}
    Under a local Gaussian background assumption, we prove that the per-pixel posterior
    probability of belonging to the target class admits a closed-form sigmoid expression:
    \begin{equation}
        W(i) = \sigma\!\left(\tau \cdot \frac{I(i) - \mu_\mathrm{ctx}}{\sigma_\mathrm{ctx}}\right),
        \label{eq:bsnr}
    \end{equation}
    where $\mu_\mathrm{ctx}$ and $\sigma_\mathrm{ctx}$ are local background statistics
    estimated from an expanded context box surrounding the annotation, and $\tau$ controls
    decision sharpness.
    This formulation bridges classical SNR analysis with Bayesian posterior estimation,
    providing a statistically grounded soft mask in a single forward pass.

    \item \textbf{Physics-Anchored Gaussian Refinement (PAG).}
    To further suppress residual clutter in the B-SNR mask, we exploit the physical
    point-spread characteristic of infrared targets.
    Instead of placing the Gaussian anchor at the \emph{geometric center} of the box
    (which is biased by annotation error), we locate the \emph{physical hotspot}
    $\mathbf{p}^* = \arg\max_i W(i)$---the pixel of maximal radiometric response.
    Using the full-width-at-half-maximum (FWHM) region of the B-SNR response as a
    clutter-free anchor, we perform SNR-weighted spatial covariance estimation to derive
    an anisotropic Gaussian prior.
    The resulting PAG mask concentrates supervision on the target's energy peak while
    gracefully attenuating toward the periphery, yielding a soft label consistent with
    the underlying point-spread function (PSF).
\end{enumerate}

Beyond pseudo-label generation, we provide a \textbf{theoretical analysis of the
variance corruption phenomenon} that arises when targets are extremely small relative
to the bounding box.
We derive a closed-form bound showing that the estimated background variance is inflated
by $\alpha(1-\alpha)\Delta\mu^2$, where $\alpha$ is the target-to-box area ratio and
$\Delta\mu$ the target--background intensity contrast.
This result rigorously explains why physics-based pseudo-labels degrade under extreme
scale conditions (sub-$3\times3$ targets in NUDT-SIRST), establishing a principled
failure boundary rather than concealing the limitation.

The main contributions of this work are summarized as follows:

\begin{itemize}
    \item We present \textbf{HALO}, the first box-supervised framework for infrared
    small target detection.
    HALO generates high-quality pixel-level pseudo-masks from coarse bounding boxes
    in a single forward pass with zero learnable parameters and no iterative refinement,
    making it a purely physics-driven, plug-and-play module compatible with any IRSTD
    backbone.

    \item We provide a \textbf{rigorous statistical foundation} for the proposed
    pseudo-label pipeline.
    B-SNR is derived as a Bayesian posterior under a Gaussian background shift model;
    PAG is formulated as SNR-weighted spatial covariance estimation anchored at the
    physical hotspot.
    We further derive a \textbf{variance corruption bound} that quantitatively
    characterizes the degradation regime under extreme target-to-box ratios.

    \item We conduct \textbf{extensive model-agnostic experiments} across three standard
    IRSTD benchmarks (NUAA-SIRST, NUDT-SIRST, IRSTD-1K) and three architecturally
    diverse backbones (DNANet, ACM, ALCNet).
    HALO consistently outperforms prior box-supervised baselines (Box Fill and GrabCut),
    surpasses single-point supervised LESPS by up to $+11.7\%$ mIoU, and achieves
    performance comparable to or exceeding PAL---all while requiring strictly coarser
    supervision and zero extra annotation effort beyond what upstream detection
    systems already produce.
\end{itemize}


% ===========================================================================
\section{Related Work}\label{sec:related}
% ===========================================================================

\subsection{Deep Learning for Infrared Small Target Detection}
\label{sec:related_dl}

The past five years have witnessed rapid architectural innovation in fully supervised
IRSTD.
DNANet~\cite{dnanet} introduced dense nested connections within a U-Net++ topology
coupled with channel--spatial attention (Res-CBAM), establishing strong segmentation
performance on NUDT-SIRST.
UIU-Net~\cite{uiunet} proposed a nested ``U-Net in U-Net'' structure to capture both
macro-level context and micro-level target details.
ACM~\cite{acm} designed an asymmetric contextual modulation module for multi-scale
feature fusion, while ALCNet~\cite{alcnet} incorporated local contrast attention into
an FPN decoder.
More recently, SCTransNet~\cite{sctransnet} and ILNet~\cite{ilnet} have explored
Transformer-based and low-level feature enhancement strategies, respectively.
Despite strong performance, all these methods critically depend on pixel-level mask
annotations, motivating the exploration of weaker supervision signals.

\subsection{Weakly Supervised Infrared Small Target Detection}
\label{sec:related_ws}

Weakly supervised IRSTD has emerged as a promising direction to reduce annotation costs.
LESPS~\cite{lesps} pioneered single-point supervision, proposing a mapping degeneration
framework in which pseudo-masks are iteratively evolved during training via the network's
own predictions.
PAL~\cite{pal} extended this paradigm with a progressive active learning framework that
strategically selects informative annotation points.

While these methods represent significant advances, they share structural limitations.
Both rely on multi-round iterative procedures that increase training complexity, and both
depend on the single-point annotation paradigm, which is fragile for extremely dim targets
(SCR $<$ 2) in heavy clutter, where clicking the target centroid demands high expertise
and is susceptible to spatial jitter.
Our work departs from this paradigm entirely, leveraging bounding box annotations---a
supervision signal naturally produced by upstream detection systems that requires no
sub-pixel human precision.

\subsection{Box-Supervised Semantic Segmentation}
\label{sec:related_box}

\paragraph{Natural image methods.}
Box-supervised segmentation has been extensively studied for natural images.
Early work by BoxSup~\cite{boxsup} and SDI~\cite{sdi} used bounding boxes to generate
coarse pseudo-masks via GrabCut~\cite{grabcut}, followed by iterative network retraining.
BBTP~\cite{bbtp} reformulated box supervision as a multiple-instance learning problem,
treating box-interior pixels as a ``bag'' with at least one foreground member.
DiscoBox~\cite{discobox} and BoxLevelSet~\cite{boxlevelset} introduced teacher--student
consistency regularization and level-set energy minimization, respectively, to sharpen
box-derived pseudo-boundaries.
BoxInst~\cite{boxinst} incorporated projection-based consistency losses within an instance
segmentation head, while Box2Mask~\cite{box2mask} proposed pairwise pixel-affinity
constraints to propagate foreground labels from box interiors.
More recent work (SIM~\cite{sim}, BoxSnake~\cite{boxsnake}) further exploits structural
image features and deformable snake contours to achieve near-mask accuracy from box
annotations alone on standard benchmarks.
A common thread across all these methods is the reliance on \emph{semantic discriminability}
or \emph{color affinity}: GrabCut requires color-based foreground--background separation;
CAM-based methods assume spatially extended class-specific activations; affinity
propagation assumes locally consistent appearance within object boundaries.

\paragraph{Cross-domain extension.}
In medical image segmentation, Kervadec~\emph{et al.}~\cite{kervadec2020weakly} derived
a differentiable loss formulation that constrains foreground size to match a box-supplied
target area estimate, avoiding the appearance assumptions of GrabCut.
Subsequent medical works (GSR-Net~\cite{gsrnet}, and~\cite{boxseg_miccai2022}) noted
explicitly that GrabCut \emph{fails on grayscale and low-contrast modalities} such as
ultrasound and MRI, motivating constraint-based or Transformer-based alternatives.
In remote sensing, \cite{boxseg_isprs2020} and SAM-based approaches~\cite{sam_rs2024}
leverage strong pre-trained features and region proposals derived from multi-spectral
cues---none of which are available in single-band thermal IR imagery.
The IR-specific method WeCoL~\cite{wecol} explores contrastive learning under weak
supervision but still requires sufficient spatial extent (targets $\geq$$6\times 6$ px).

\paragraph{Failure modes on infrared small targets.}
Despite extensive progress across domains, \emph{all} of the above methods share a
critical prerequisite that is systematically violated in IRSTD: the target must be
distinguishable from its local background by appearance, semantics, or learned
features.
An infrared point target typically occupies fewer than $9\times9$ pixels---frequently
as few as $1\times1$---and presents as a near-isotropic intensity blob with no texture,
no color, and no semantic structure.
Its surrounding background clutter has a statistically similar intensity distribution,
making GrabCut's Gaussian mixture model degenerate, CAM responses indistinguishable
from noise, and pairwise affinity constraints unreliable.
This failure is not merely empirical: the MICCAI 2022 medical segmentation
study~\cite{boxseg_miccai2022} identifies the same mechanism when GrabCut is applied to
low-contrast grayscale modalities, confirming that the breakdown is \emph{modality-driven},
not dataset-specific.
The closest IR-native approaches, LESPS~\cite{lesps} and HMG~\cite{hmg}, circumvent
this limitation by replacing appearance-based cues with network-iterative pseudo-mask
feedback, at the cost of multi-round training dependencies.
Our work takes an orthogonal route: rather than improving the discriminability of
learned features, we exploit the known \emph{radiometric physics} of infrared point
sources---local SNR and the optical point-spread function---to derive pseudo-masks
offline, in a single pass, with zero learnable parameters.


% ===========================================================================
\section{Methodology}\label{sec:method}
% ===========================================================================

\subsection{Problem Formulation}
\label{sec:method_prob}

Let $\mathbf{I} \in \mathbb{R}^{H \times W}$ denote a grayscale infrared image.
We assume access to a set of $K$ bounding box annotations
$\mathcal{B} = \{(x_k^1, y_k^1, x_k^2, y_k^2)\}_{k=1}^K$ obtained from an upstream
detection system (e.g., a YOLO detector or radar cueing).
Each box is a loose enclosure of a single infrared small target; no pixel-level mask
or target centroid is available.
Our goal is to generate a pixel-level pseudo-label map
$\hat{\mathbf{Y}} \in [0,1]^{H \times W}$ that can directly supervise any IRSTD
segmentation network $f_\theta$, without introducing learnable parameters or iterative
training loops.

The generation pipeline, illustrated in Fig.~\ref{fig:pipeline}, proceeds in two
physically motivated stages: (1) B-SNR posterior estimation, which provides
a per-pixel foreground probability within each box, and (2) PAG refinement, which
further suppresses residual clutter by anchoring a Gaussian prior at the physical
radiometric hotspot.
Both stages rely exclusively on the local statistics of the input image and require
no data-driven optimization.

\subsection{Stage 1: B-SNR Posterior Mask}
\label{sec:method_bsnr}

\paragraph{Statistical background model.}
Infrared background clutter at local scales is well-approximated by a Gaussian
distribution~\cite{irstd_survey}.
For a target box $b = (x^1, y^1, x^2, y^2)$, we define an expanded \emph{context box}
$b_\mathrm{ctx}$ by enlarging $b$ by a factor of $\rho_\mathrm{ctx} = 1.5$ in each
spatial dimension.
The context region $\Omega_\mathrm{ctx} = b_\mathrm{ctx} \setminus b$ (annular region
surrounding the target box) provides background pixels free from target contamination.
We estimate the local background statistics as:
\begin{equation}
    \mu_\mathrm{ctx} = \frac{1}{|\Omega_\mathrm{ctx}|} \sum_{i \in \Omega_\mathrm{ctx}} I(i),
    \qquad
    \sigma_\mathrm{ctx}^2 = \frac{1}{|\Omega_\mathrm{ctx}|} \sum_{i \in \Omega_\mathrm{ctx}}
    \bigl(I(i) - \mu_\mathrm{ctx}\bigr)^2.
    \label{eq:bg_stats}
\end{equation}

\paragraph{Bayesian posterior derivation.}
Under a binary hypothesis test, pixel $i$ inside box $b$ belongs either to the
background class ($y=0$, modeled as $I(i) \sim \mathcal{N}(\mu_\mathrm{ctx},
\sigma_\mathrm{ctx}^2)$) or to the target class ($y=1$, modeled as
$I(i) \sim \mathcal{N}(\mu_\mathrm{ctx} + \Delta\mu, \sigma_\mathrm{ctx}^2)$ with
$\Delta\mu > 0$).
Applying Bayes' theorem with equal priors, the posterior probability of foreground is:
\begin{align}
    P(y{=}1 \mid I(i)) &= \sigma\!\left(\frac{\log P(I(i)|y{=}1) - \log P(I(i)|y{=}0)}{1}\right) \notag \\
    &= \sigma\!\left(\frac{2\Delta\mu}{\sigma_\mathrm{ctx}^2}\cdot I(i)
       - \frac{\Delta\mu\,(2\mu_\mathrm{ctx}+\Delta\mu)}{\sigma_\mathrm{ctx}^2}\right).
    \label{eq:posterior_full}
\end{align}
Since $\Delta\mu$ is unknown a priori, we treat the coefficient as a single
temperature parameter $\tau > 0$ and center the response at the estimated background
mean, yielding the \textbf{B-SNR weight}:
\begin{equation}
    \boxed{W(i) = \sigma\!\left(\tau \cdot \frac{I(i) - \mu_\mathrm{ctx}}
    {\sigma_\mathrm{ctx} + \varepsilon}\right)}, \quad i \in b,
    \label{eq:bsnr_final}
\end{equation}
where $\varepsilon = 10^{-4}$ prevents division by zero.
Outside box $b$, $W(i) = 0$.
The SNR ratio $(I(i)-\mu_\mathrm{ctx})/\sigma_\mathrm{ctx}$ is the classical
signal-to-clutter ratio (SCR) used in traditional IRSTD detectors~\cite{irstd_survey};
our formulation gives it a Bayesian probabilistic interpretation.
After computing $W$, we apply per-box min-max normalization to map the response to
$[0, 1]$, ensuring that the peak foreground probability is always 1 regardless of
absolute intensity.

\subsection{Stage 2: Physics-Anchored Gaussian Refinement (PAG)}
\label{sec:method_pag}

\paragraph{Motivation.}
The B-SNR mask recovers high-response pixels within the box but may retain residual
clutter responses away from the target peak, particularly when background statistics
are locally non-stationary.
To further focus supervision on the true target locus, we exploit the \emph{point-spread
function} (PSF) of infrared optical systems: the energy distribution of a distant
point source follows a Gaussian-like profile whose spread is determined by optical
diffraction and atmospheric turbulence.

\paragraph{Physical hotspot localization.}
Rather than using the \emph{geometric center} of the bounding box as the Gaussian
anchor---which is biased by annotation error and does not reflect the true target
location---we define the \textbf{physical hotspot} as:
\begin{equation}
    \mathbf{p}^* = \arg\max_{i \in b}\, \mathrm{SNR}(i),
    \qquad
    \mathrm{SNR}(i) = \frac{I(i) - \mu_\mathrm{ctx}}{\sigma_\mathrm{ctx} + \varepsilon}.
    \label{eq:hotspot}
\end{equation}
Because the SNR is maximized precisely at the pixel of highest radiometric response,
$\mathbf{p}^*$ is insensitive to box boundary perturbations as long as the target
is enclosed---a key source of HALO's robustness to annotation noise, as confirmed
in our perturbation experiments (Sec.~\ref{sec:exp_robustness}).

\paragraph{Adaptive Gaussian bandwidth estimation.}
We estimate the spatial spread of the target response using only the dominant energy
peak, excluding secondary clutter spikes.
Define the FWHM region as:
\begin{equation}
    \Omega_\mathrm{FWHM} = \bigl\{i \in b : \mathrm{SNR}(i) \geq
    {\textstyle\frac{1}{2}}\,\mathrm{SNR}(\mathbf{p}^*)\bigr\}.
    \label{eq:fwhm}
\end{equation}
Within $\Omega_\mathrm{FWHM}$, we compute the SNR-weighted mean-square distance to
the hotspot as an estimate of the target's spatial covariance:
\begin{equation}
    \hat{\sigma}^2 = \gamma^2 \cdot
    \frac{\sum_{i \in \Omega_\mathrm{FWHM}} \mathrm{SNR}^+(i)\,\|\mathbf{r}_i\|^2}
         {\sum_{i \in \Omega_\mathrm{FWHM}} \mathrm{SNR}^+(i)},
    \label{eq:pag_sigma}
\end{equation}
where $\mathrm{SNR}^+(i) = \max(\mathrm{SNR}(i), 0)$,
$\mathbf{r}_i = \mathbf{p}_i - \mathbf{p}^*$ is the displacement from the hotspot,
and $\gamma = 1.5$ is a conservative scale factor ensuring the Gaussian envelope
covers the target boundary.
A lower bound of $\hat{\sigma} \geq 1.0$ pixel prevents degenerate collapse on
sub-pixel targets.

\paragraph{PAG soft label.}
The final PAG pseudo-label map is computed as the product of the B-SNR weight and
the physics-anchored Gaussian:
\begin{equation}
    \boxed{P(i) = W(i) \cdot \exp\!\left(-\frac{\|\mathbf{r}_i\|^2}{2\hat{\sigma}^2}\right)},
    \quad i \in b,
    \label{eq:pag_final}
\end{equation}
and $P(i) = 0$ outside $b$.
The B-SNR factor $W(i)$ provides background suppression from local statistics, while
the Gaussian factor concentrates the pseudo-label near the physical hotspot.
Their product yields a soft label that simultaneously respects both the statistical
SNR evidence and the optical PSF geometry.
When multiple targets are present, PAG is applied per box and the per-pixel maximum
is taken across all targets.

\subsection{Variance Corruption Bound}
\label{sec:method_theory}

We now analyze theoretically when and why the B-SNR estimation degrades.
Consider a box $b$ of area $S$ enclosing a compact target of area $s$, so that the
target-to-box area ratio is $\alpha = s/S$.
Let background pixels have intensity distribution $\mathcal{N}(\mu_0, \sigma_0^2)$ and
target pixels have mean $\mu_\mathrm{tgt} = \mu_0 + \Delta\mu$ with $\Delta\mu > 0$.

In practice, the context statistics $\mu_\mathrm{ctx}$ and $\sigma_\mathrm{ctx}^2$ are
estimated over the expanded context box $b_\mathrm{ctx}$, which contains a small
fraction $\alpha'$ of target pixels.
By the law of total variance applied to the mixture distribution in $b_\mathrm{ctx}$:

\begin{proposition}[Variance Corruption Bound]
\label{prop:vcb}
If the target-to-context-box area ratio is $\alpha' \approx \alpha / \rho_\mathrm{ctx}^2$,
the estimated context variance satisfies:
\begin{equation}
    \hat{\sigma}_\mathrm{ctx}^2 \approx \sigma_0^2
    + \underbrace{\alpha'(1 - \alpha')\Delta\mu^2}_{\text{corruption term}}.
    \label{eq:vcb}
\end{equation}
\end{proposition}

\begin{proof}
The mixture variance of a two-component population with proportions $(1-\alpha', \alpha')$
and means $(\mu_0, \mu_0 + \Delta\mu)$ (variance $\sigma_0^2$ each) equals:
$\hat{\sigma}^2 = (1-\alpha')\sigma_0^2 + \alpha'\sigma_0^2
+ (1-\alpha')\alpha'^2\Delta\mu^2 + \alpha'(1-\alpha')^2\Delta\mu^2
= \sigma_0^2 + \alpha'(1-\alpha')\Delta\mu^2$.
\hfill$\square$
\end{proof}

\paragraph{Implications.}
The corruption term $\alpha'(1-\alpha')\Delta\mu^2$ is negligible when $\alpha' \to 0$
(target much smaller than context box, $\hat{\sigma}_\mathrm{ctx} \approx \sigma_0$,
B-SNR accurate) or $\alpha' \to 1$ (box tightly surrounds target, rarely occurs).
The maximum corruption occurs at $\alpha' = 0.5$.
For the NUDT-SIRST dataset, targets frequently occupy sub-$3\times3$ pixels with
high contrast ($\Delta\mu \gg \sigma_0$): even though $\alpha' \ll 1$, the term
$\alpha'\Delta\mu^2$ can dominate $\sigma_0^2$, causing $\hat{\sigma}_\mathrm{ctx}$
to be severely inflated and driving the B-SNR response of the target toward $0.5$
rather than $1$.
This analytically explains the performance gap between NUDT-SIRST and NUAA-SIRST
observed in our experiments---the limitation is not a design flaw but a theoretically
characterized failure mode inherent to any SNR-based method operating near the
physical detectability limit.

\subsection{Integration with IRSTD Networks}
\label{sec:method_train}

HALO is strictly a \emph{pseudo-label generator}; it introduces no changes to network
architecture, training loss, or inference procedure.
Given a training image $\mathbf{I}$ with box annotations $\mathcal{B}$, we pre-compute
the PAG pseudo-label map $\hat{\mathbf{Y}}$ offline before training begins.
The resulting soft mask is stored alongside the image and fed to the segmentation network
in place of the ground-truth mask.

\paragraph{Training objective.}
We adopt the Soft-IoU loss~\cite{softiou}, which operates directly on continuous
predictions and soft labels without requiring binarization:
\begin{equation}
    \mathcal{L}_\mathrm{SoftIoU}(f_\theta(\mathbf{I}),\,\hat{\mathbf{Y}})
    = 1 - \frac{\sum_i f_\theta(\mathbf{I})_i \cdot \hat{Y}_i}
               {\sum_i \bigl(f_\theta(\mathbf{I})_i + \hat{Y}_i
               - f_\theta(\mathbf{I})_i \cdot \hat{Y}_i\bigr)},
    \label{eq:softiou}
\end{equation}
where $f_\theta(\mathbf{I})_i \in [0,1]$ is the sigmoid-activated network output
at pixel $i$.
The soft label $\hat{Y}_i \in [0,1]$ encodes both the foreground/background decision
and the confidence gradient, allowing the network to learn a smooth activation
concentrated at the physical hotspot rather than a hard binary mask.

\paragraph{Offline pre-generation.}
Since HALO requires no information from the network (it is computed solely from $\mathbf{I}$
and $\mathcal{B}$), pseudo-labels can be pre-generated once and cached as PNG files
before training.
This eliminates the per-iteration overhead of iterative methods such as LESPS and
introduces zero additional GPU memory cost at training time.
The entire pseudo-label dataset for all three benchmarks can be generated in under
two minutes on a standard CPU.

\paragraph{Model-agnostic plug-in.}
Because HALO only replaces the mask annotation file, it is trivially compatible with
any IRSTD network that accepts pixel-level supervision.
We validate this property across three architecturally diverse backbones (DNANet,
ACM, ALCNet) in Sec.~\ref{sec:exp_agnostic}, observing consistent and significant
gains over the Box Fill baseline, which confirms that the performance improvement
originates from the physics-anchored pseudo-label quality rather than any
backbone-specific property.
